\chapter{Що машина робить уночі}
% full version: chapters/20_sleep.tex

Сон --- не вимкнення машини, а зміна режиму. Це видно з радіостанцій. У глибокому сні приглушені майже всі: ацетилхолін, норадреналін, серотонін. У сні зі сновидіннями ацетилхолін повертається до денного рівня, норадреналін і серотонін майже мовчать, а вихід на м'язи вимкнено гальмами --- тому ми не тікаємо насправді від того, від чого тікаємо уві сні. Усе це виміряно.

\section*{Коли час спати}

Уявіть конденсатор, який заряджається весь день, поки ви не спите, і розряджається вночі. Заряд --- це «тиск сну». Засинаєте ви, коли заряд сягає верхнього порогу, прокидаєтеся, коли він опускається до нижнього. А самі пороги плавно ходять угору й униз разом із добовим ритмом із розділу 16. Така проста схема сама приходить до шістнадцяти годин неспання й восьми годин сну. І пояснює, чому після безсонної ночі ми спимо не вдвічі довше, а глибше: повний конденсатор розряджається швидше. Кандидат на роль «заряду» --- речовина, що накопичується в мозку під час роботи; що це саме вона, поки гіпотеза.

\pencilfig{120}{%
% figures/sleep_{S,H,L}.dat from models/sleep_models.py, t in hours
\begin{scope}[x=0.1cm, y=2.6cm]
\foreach \a/\b in {0/4.3,21.2/29.3,45.4/53.4,69.5/72}{\fill[pcBlue, opacity=0.1] (\a,0) rectangle (\b,1);}
\draw[pen] (0,0) -- (72,0);
\draw[penlite=pcRed, densely dashed] plot file {figures/sleep_H.dat};
\draw[penlite=pcGreen, densely dashed] plot file {figures/sleep_L.dat};
\draw[pcBlue, line width=1.0pt] plot file {figures/sleep_S.dat};
\foreach \t/\l in {0/0,24/доба,48/дві,72/три}{\draw[pcGray, line width=0.5pt] (\t,0) -- (\t,-0.04); \node[lbl, anchor=base] at (\t,-0.16) {\l};}
\node[lbl, text=pcBlue] at (25.25,0.93) {сон}; \node[lbl, text=pcBlue] at (49.4,0.93) {сон};
\end{scope}
\node[lbl, text=pcRed, anchor=west] at (7.3,2.0) {заснути};
\node[lbl, text=pcGreen!70!black, anchor=west] at (7.3,0.62) {прокинутися};
\node[lbl, text=pcBlue, anchor=west] at (7.3,1.35) {тиск сну};
}{Тиск сну накопичується вдень і тане вночі; засинаємо біля верхнього порогу, прокидаємося біля нижнього, а обидва пороги гойдаються разом із добою.}

\section*{Повільна гойдалка}

У глибокому сні цілі ділянки мозку синхронно то вмикаються, то вимикаються --- рідше ніж раз на секунду. Це вже знайома схема: група, яка збуджує сама себе і втомлюється, --- генератор. Удень ацетилхолін пригнічує втому нейронів, і та сама група працює рівно; уночі його мало, і мережа починає гойдатися.

\pencilfig{20}{\pencilart{20}}{Уночі мережа гойдається повільно: то майже всі нейрони працюють, то майже всі мовчать.}

\section*{Перемотування вперед}

У цій гойдалці відбувається найцікавіше. Ланцюжки спрацювань нейронів, що сталися вдень, «програються» знову --- і в рази швидше: у відділі пам'яті приблизно вдвадцятеро, у корі в кілька разів. Якщо штучно посилити узгодженість цих програвань із повільною гойдалкою, пам'ять поліпшується --- це причиновий дослід. Навіщо це потрібно? Правдоподібна гіпотеза: повільна довготривала пам'ять учиться на повторах упереміш зі старим, щоб нове не затерло попереднього. Інженери знають цю біду: нейромережа, навчена на новій задачі, забуває стару, якщо не повторювати старих прикладів.

\section*{Підрізання зв'язків}

Ще одна виміряна річ: після сну контакти між нейронами в середньому стають меншими --- на одиниці-десятки відсотків, пропорційно до свого розміру. Ніби всі гучності трохи зменшили одним рухом: співвідношення збереглися, вивчене лишилося, а місце для завтрашнього навчання звільнилося. Що це головне завдання сну, --- гіпотеза.

\section*{Сновидіння і прибирання}

Сон зі сновидіннями схожий на стендове випробування: машина працює на повну потужність, але входи замінено внутрішніми сигналами, а вихід вимкнено. Навіщо це потрібно, ми не знаємо. Ідея, що уві сні мозок «промивається» від відходів, поки спірна: є досліди і за, і проти.

Найдивовижніше: сон буває місцевим. Після довгого неспання окремі ділянки кори людини, яка не спить, на мить «засинають», замовкаючи посеред роботи.

\fullversion{20}
